\chapter{Procedures the chapters call for}\label{app:procedures}

Several chapters say that an input matters and stop short of saying how to set it. This
appendix gives the procedure for each: arguing for a reading precision, dividing the error
budget, choosing a remedy when programme risk exceeds the tolerance, defending a tolerance
to an authority, repairing a failed test, and weighing who bears the harm. The arithmetic is
in \texttt{privacy\_risk.procedures} and \texttt{privacy\_risk.validation}, with tests in
\texttt{tests/test\_ext\_procedures.py}. The procedures organise an institution's judgement.
They do not supply the inputs that only the institution holds.

\section{Arguing for a reading precision}\label{app:b-precision}

The adversary reads a continuous variable, here log exposure, with an error of standard
deviation $s$. Chapter~\ref{ch:credit} declares $s = 0.6$ and reports floors at 0.3 and 1.0.
The institution has to say why its own value is not very different.

\begin{enumerate}
\item \textbf{List the sources an outsider can hold.} Commercial databases, public filings,
the outsider's own transactions with the borrower, and published accounts are the usual
ones. Name each one and who can buy or read it.
\item \textbf{Bound the error of each source.} A figure rounded to the nearest $h$ carries a
rounding error of standard deviation $h/\sqrt{12}$, and on the log scale that is divided by
the level. Exposure rounded to the nearest 1,000 on a loan of 50,000 gives
$1000/\sqrt{12}/50000 = 0.0058$. A figure that is stale, estimated or modelled needs an
error taken from a comparison with the institution's own records, not from the source's
description of itself.
\item \textbf{Combine sources only if their errors are independent.} Independent unbiased
readings combine to a standard deviation of $\bigl(\sum_k s_k^{-2}\bigr)^{-1/2}$, which is
below the best single source: two readings of 0.6 give $0.6/\sqrt{2} = 0.424$, and a source
of 0.3 joined by one of 1.0 gives 0.287. Sources that copy one another add nothing, so the
institution should say what it knows about the lineage of each.
\item \textbf{Declare the smallest value the evidence allows, then report the floors across a
range.} The decision uses the declared precision. The record also shows the floor at a
coarser and a finer value, so that a reader can see how far the answer depends on the
argument.
\end{enumerate}

Table~\ref{app:tab-prec} shows step four on the repository's synthetic register, with the
fixed-sequence procedure of Chapter~\ref{ch:selection} over the floors 150, 75, 50, 20, 10, 5
and 1, a tolerance of 0.12, $\alpha = 0.05$ and 1,500 attacks per floor. The attack is the
repository's single-release linkage on the borrower key, which is not the attack that
Chapter~\ref{ch:credit} uses to select its floors, so these floors do not reproduce that
chapter's and are not meant to. What the table shows is the shape of the dependence.

\begin{table}[t]
\centering
\caption{Exact upper bounds at each floor and the floor the fixed sequence selects, by
reading precision. Synthetic register of 32,581 loans; 1,500 attacks per cell; seed from
\texttt{configs/paper.json}. A dash means the sequence had stopped.}\label{app:tab-prec}
\small
\begin{tabular}{@{}lrrrrrrrl@{}}
\toprule
$s$ & 150 & 75 & 50 & 20 & 10 & 5 & 1 & Selected \\
\midrule
0.3 & 0.035 & 0.052 & 0.062 & 0.107 & 0.138 & --- & --- & 20 \\
0.6 & 0.024 & 0.031 & 0.046 & 0.067 & 0.086 & 0.088 & 0.088 & none \\
1.0 & 0.010 & 0.022 & 0.029 & 0.039 & 0.055 & 0.065 & 0.070 & none \\
\bottomrule
\end{tabular}

\end{table}

A sharper adversary needs a floor of 20 where a coarser one needs none, so the precision is
a decision input. An institution that cannot argue for 0.6 over 0.3 should adopt the floor
that 0.3 requires.

\section{Setting the error budget}\label{app:b-budget}

Chapter~\ref{ch:monitoring} gives two rules: the cost rule $w_i \propto c_i/\mu_i$ for
coordinates known in advance, and the schedule $w_i = 6/(\pi^2 i^2)$ for an open-ended
family. The procedure joins them.

\begin{enumerate}
\item Fix the total error $\alpha$ for the programme, for example 0.05. It does not change
once outcomes have been seen.
\item Decide a reserve fraction $\rho$ for coordinates not yet proposed. The reserve costs the
known coordinates a factor $1-\rho$ of their budget, which is the trade-off to state: a large
reserve makes every audit under the budget longer, and a small one leaves no room for a
stronger adversary.
\item Divide $(1-\rho)\alpha$ among the known coordinates in proportion to $c_i/\mu_i$, with
the costs and planning drifts fixed before any outcome.
\item Spread $\rho\alpha$ over later coordinates by $6/(\pi^2 j^2)$.
\item When a product is added, pay for it from the reserve. If the reserve cannot cover it,
the total cannot be raised afterwards, and the choice is a smaller share or a later
programme.
\end{enumerate}

For $\alpha = 0.05$, $\rho = 0.2$ and three known attacks with costs 1, 2 and 4 per draw and
planning drifts 0.02, 0.02 and 0.04, the weights $c_i/\mu_i$ are 50, 100 and 100, and the
shares are 0.008, 0.016 and 0.016. The log-wealth thresholds $\ln(1/w_i)$ are 4.83, 4.14 and
4.14. The reserve is 0.010, and the first five later coordinates receive 0.00608, 0.00152,
0.00068, 0.00038 and 0.00024.

\section{Choosing a remedy for the composition gap}\label{app:b-remedy}

Chapter~\ref{ch:programmes} shows that no function of per-release risks bounds the pooled
risk, so the remedy is chosen by re-simulating the pooled attack, not by arithmetic on the
releases. The tree has three branches.

\begin{description}
\item[Small gap.] The upper bound on the pooled risk is within the programme tolerance.
Release as planned.
\item[Moderate gap.] The pooled bound exceeds the tolerance, but some remedy restores it. Try
the remedies in order of cost, coarsening linking fields, rotating pseudonyms, reducing
frequency, then moving to a secure environment, and re-run the pooled attack after each. Take
the first that brings the bound within the tolerance. A remedy is not credited until the
pooled attack has been run against it.
\item[Large gap.] No remedy restores the tolerance. Change the product, for example to
aggregates, and assess the new product from the start.
\end{description}

The rule is \texttt{procedures.\allowbreak{}composition\_\allowbreak{}remedy}. In Chapter~\ref{ch:falsification}
the declared four-release programme has a pooled bound of 0.464 against a tolerance of 0.12, so
it sits on the second or third branch, and which one depends on re-runs that the repository
has not made. The inputs in the function's test are illustrative.

\section{Defending the tolerance}\label{app:b-legal}

A data protection authority may ask why a tolerance of 0.12 is a reasonable reading of
``reasonably likely''. The record answers with items that can be checked for presence, though
not for correctness. The institution should hold:

\begin{enumerate}
\item the legal reading, written down and dated, with the provision it interprets;
\item the reasons for the number, in terms of the harm class, the protected parties and the
access route;
\item the owner, who is senior enough to answer for it in public;
\item the trigger for review: a change in the law, the technology or the access route, and the
expiry date.
\end{enumerate}

\texttt{procedures.DefenceRecord} lists these and reports which are blank. It checks
completeness. Whether the reading is reasonable is for lawyers, and the book does not claim to
settle it.

\section{Repairing a failed test}\label{app:b-repair}

A failed falsification test makes the assessment invalid. The repair keeps the failure on the
record.

\begin{enumerate}
\item Record the failed result, the band and the data it was run on. It is not overwritten.
\item Name the assumption the test touches. A calibration failure points at the error scale or
the linkage model, a transport failure at the frequency model, a monotonicity failure at the
mechanism or the noise.
\item Revise the model on development data only.
\item Fix the new band before the rerun, from the size of error that would change the
decision, and say so in the record.
\item Rerun on fresh draws. The earlier run motivated the band, so it cannot be used to show
the band valid.
\end{enumerate}

The calibration band of 0.03 and the uniqueness limit of 0.30 are both at step three of this
list for that reason: each was set after the evidence was seen.

\section{Who bears the harm}\label{app:b-ethics}

The method treats harm as a probability inside a declared harm class, and a tolerance as a
social choice. Four questions sit outside what it computes, and an institution should answer
them in the record.

\paragraph{Whose interests.} The institution wants defensible decisions, data subjects do not
want to be identified, and users want useful data. The tolerance and the floor are where those
interests are weighed, and the owner should say how.
\paragraph{Distribution of harm.} A risk of 0.12 averaged over targets can hide a group for
whom it is far higher. The worst-target and subgroup bounds of Chapters~\ref{ch:estimate} and
\ref{ch:credit} are the tools; the record should name the groups examined.
\paragraph{Consent.} The legal basis may not need the subjects' consent. Whether it should be
sought is a question the method cannot answer.
\paragraph{Accountability.} The owner accepts the residual risk. The record should say what
follows if a release identifies a person: who is told, who answers for it, and what is done
about the programme.
